% =====================================================================
% Clase -- Cálculo 11° -- Trimestre I -- Semana 01
% Sesiones 001 (lun 3 ago., 100 min) y 002 (jue 6 ago., 55 min)
% Tema 1 de la guía: Los sistemas numéricos
% Material del docente (no se entrega a los estudiantes).
% =====================================================================
\documentclass[11pt]{article}
\makeatletter\def\input@path{{../../../../../plantillas/estilo/}}\makeatother
\usepackage{mmcantillo}
\guiadelestudiante{../../guia-didactica/guia-periodo-I-numeros-reales}

\tipodocumento{Clase}
\asignatura{Cálculo}
\titulo{Semana 1: los sistemas numéricos}
\grado{11°}
\periodo{I}
% DBA: matematicas grado 11 · DBA 1 — Utiliza las propiedades de los números (naturales,
%      enteros, racionales y reales) y sus relaciones y operaciones para construir y comparar
%      los distintos sistemas numéricos.
% Evidencias trabajadas: (1) describe propiedades de los números y las operaciones que son
%      comunes y diferentes en los distintos sistemas numéricos.

\begin{document}

\encabezado*

\begin{center}
{\renewcommand{\arraystretch}{1.3}
\begin{tabularx}{\linewidth}{@{}l l l X l@{}}
  \toprule
  \textbf{Sesión} & \textbf{Fecha} & \textbf{Min} & \textbf{Subtema} & \textbf{Entregables}\\
  \midrule
  001 & lun 3 ago., 07:50 & 100 & Diagnóstico; números naturales y enteros & Taller 1\\
  002 & jue 6 ago., 06:55 & 55  & Números racionales & Tarea 1 (para el lun 10 ago.)\\
  \bottomrule
\end{tabularx}}
\end{center}

\small El viernes 7 de agosto es festivo (Batalla de Boyacá): no hay clase. \normalsize
Referencia: \refguia{tema:sistemas}.

% ---------------------------------------------------------------------
\seccion{Sesión 001 — lunes 3 de agosto · 07:50 · 100 min}

\textbf{Objetivo.} Presentar el curso y la guía del trimestre, diagnosticar lo que el grupo
recuerda de los números, y reconstruir por qué se pasa de $\N$ a $\Z$.

\textbf{Materiales.} Guía del estudiante impresa; Taller 1 (una copia por estudiante); tablero.

\begin{center}
{\renewcommand{\arraystretch}{1.3}
\begin{tabularx}{\linewidth}{@{}l l X@{}}
  \toprule
  \textbf{Momento} & \textbf{Min} & \textbf{Actividad}\\
  \midrule
  Presentación & 0–10 & El curso en una frase: «el cálculo es la matemática del cambio,
    y descansa sobre la recta real». Escala de notas de 1,0 a 5,0; cómo funcionan
    guía, talleres, quices y tareas.\\
  Diagnóstico & 10–30 & Cinco preguntas en el tablero, individuales y sin nota (ver abajo).
    Se corrigen en voz alta; anota quiénes fallan en signos y fracciones.\\
  Hilo & 30–45 & Lectura en voz alta de la introducción de la guía y del hilo del Tema 1
    (los pitagóricos: «todo es número»). Pregunta: ¿qué significa que la música
    «sea» números?\\
  Explicación & 45–65 & $\N$ sirve para contar; la resta $3 - 5$ no tiene respuesta en $\N$
    y por eso aparece $\Z$. Recta numérica con enteros; opuesto de un número;
    $\N \subset \Z$. Operación \emph{cerrada}: su resultado siempre queda en el conjunto.\\
  Taller 1 & 65–95 & En parejas. Circula por los puestos; en el punto 3, quien diga que el
    punto medio entre $-4$ y $5$ es un entero tiene la pregunta de la próxima clase.\\
  Cierre & 95–100 & Cada estudiante escribe en una línea: «¿Qué operación obligó a inventar
    los enteros?» Se recoge el taller.\\
  \bottomrule
\end{tabularx}}
\end{center}

\textbf{Diagnóstico (en el tablero).}
\begin{enumerate}[label=\alph*)]
  \item $-7 + 12$ \hfill \emph{Respuesta: $5$}
  \item $(-3)(-4)$ \hfill \emph{$12$}
  \item $\frac{2}{3} + \frac{1}{4}$ \hfill \emph{$\frac{11}{12}$}
  \item ¿Es cierto que $0{,}5 = \frac{1}{2}$? \hfill \emph{Sí}
  \item ¿$\sqrt{9}$ es un número natural? ¿Y $\sqrt{2}$? \hfill \emph{Sí ($3$); $\sqrt{2}$ no: se verá en el Tema 2}
\end{enumerate}

\textbf{Preguntas para el grupo.}
\begin{itemize}
  \item ¿Hay algún número natural que no tenga siguiente? ¿Y alguno que no tenga anterior?
        \emph{(Todo natural tiene siguiente; el 0 no tiene anterior en $\N$.)}
  \item Si la temperatura es $-2$ °C y baja 5 grados, ¿qué operación hacemos?
        \emph{($-2 - 5 = -7$.)}
\end{itemize}

\begin{nota}
Si el diagnóstico muestra muchos errores con signos, dedica 5 minutos de la sesión 002 a la
ley de los signos antes de los racionales. Para quienes terminen antes el taller: ¿es cerrada
la división en $\Z$? (No: $3 \div 5$.)
\end{nota}

% ---------------------------------------------------------------------
\seccion{Sesión 002 — jueves 6 de agosto · 06:55 · 55 min}

\textbf{Objetivo.} Definir los números racionales, relacionarlos con los decimales finitos y
periódicos, y convertir un decimal periódico en fracción.

\begin{center}
{\renewcommand{\arraystretch}{1.3}
\begin{tabularx}{\linewidth}{@{}l l X@{}}
  \toprule
  \textbf{Momento} & \textbf{Min} & \textbf{Actividad}\\
  \midrule
  Enganche & 0–5 & Retoma la pregunta del taller: el punto medio entre $-4$ y $5$ es
    $0{,}5$, que no es entero. «Necesitamos otro conjunto.»\\
  Explicación & 5–25 & $\Q = \{\frac{p}{q} : p, q \in \Z,\ q \neq 0\}$. División larga de
    $\frac{3}{8}$ (se acaba: $0{,}375$) y de $\frac{1}{7}$ (los restos se repiten:
    $0{,}\overline{142857}$). Idea clave: al dividir entre $q$ solo hay $q - 1$ restos
    posibles distintos de cero, así que el decimal \emph{tiene} que terminar o repetirse.\\
  Ejemplo guiado & 25–40 & El decimal periódico $0{,}\overline{12}$ como fracción (ejemplo de
    la guía): $100x - x = 12$, $x = \frac{12}{99} = \frac{4}{33}$. Luego $0{,}\overline{5}$
    ($\frac{5}{9}$) en el tablero con participación del grupo.\\
  Práctica & 40–50 & Ejercicios 1 y 2 del Tema 1 de la guía, en parejas.\\
  Cierre & 50–55 & Resumen: «racional = decimal finito o periódico». Se asigna la Tarea 1
    para el lunes 10 de agosto (el viernes 7 es festivo).\\
  \bottomrule
\end{tabularx}}
\end{center}

\textbf{Respuestas de los ejercicios de la guía} (Tema 1): ejercicio 1: $7 - 12 = -5$ en
$\Z$; $\frac{-8}{4} = -2$ en $\Z$; $\frac{2}{3} + \frac{1}{6} = \frac{5}{6}$ en $\Q$;
$5 \cdot 9 = 45$ en $\N$. Ejercicio 2: $0{,}\overline{3} = \frac{1}{3}$;
$0{,}\overline{45} = \frac{45}{99} = \frac{5}{11}$; $1{,}2\overline{5} =
\frac{125 - 12}{90} = \frac{113}{90}$.

\textbf{Preguntas para el grupo.}
\begin{itemize}
  \item ¿$0{,}999\ldots$ es menor que 1? \emph{(Son iguales: $10x - x = 9$, $x = 1$. Suele
        generar discusión: úsala para mostrar que la notación decimal no es única.)}
  \item ¿Todo número entero es racional? \emph{(Sí: $n = \frac{n}{1}$, así que $\Z \subset \Q$.)}
\end{itemize}

\end{document}
